\section{Physical scales and stopping levels}\label{sec:rg-flow}

The logarithmic charge scale determines where the finite-volume
iteration stops. We express the stopping rule in terms of depth below
the physical mass scale. This removes the cutoff from the terminal
physical quantities and gives the side margins needed for comparison
at every preceding step.
We follow \cite[\S16.1]{QED3Master}, which retains the stopping-scale
calculation of \cite[Chapter~19]{LatticeCurvature}.

\subsection{Physical scales and stopping levels}
\label{subsec:rgfl-scales}

Fix $\bar m\in(0,1]$ and put
$h=L^{\lfloor\log_L(1/\bar m)\rfloor}$.
For a cutoff $N\ge0$, define
\begin{equation}\label{eq:rgfl-physical-scales}
 \eta=L^{-N},\qquad K=N+\log_Lh,\qquad
 a_j=L^{j-N},\qquad e_j=e\sqrt{a_j},\qquad
 \mu_j=\bar m a_j.
\end{equation}
Thus $a_j$ is the physical length of one unit at
level $j$, and $a_K=h$.
At zero the charge is written $q=e\sqrt\eta$;
the fixed charge ceilings $\bar e_0,\bar e_1$ retain their
earlier meanings. A formula with variable index $j$
uses $e_j$ for the running charge, including $q$
when $j=0$.
For $0\le j\le K$ set
\[
 v_j=-\log e_j,\qquad
 r_j=\min\{L^i:i\in\N\cup\{0\},\ L^i\ge v_j^r\},
 \qquad p_j=v_j^p,\quad p_{0,j}=v_j^{p_0}.
\]
The integers $r,p,p_0$ and the constants of the
preceding estimates are fixed.
Here $r\ge2$, $L$ is odd with $L\ge13$, and $M$
is a power of $L$ with $M\ge L^2$.
We assume the hypotheses of the primitive and
step estimates, and use a ceiling
$\bar e_1\le\min(\bar e_0,1/3)$ such that, for all
$v\ge-\log\bar e_1$,
\begin{equation}\label{eq:rgfl-scale-smallness}
 \begin{gathered}
 v\ge8,\qquad
 \left(1+\frac{3\log L}{2v}\right)^r<L,\qquad
 v^r>31L+16,\\
 (c_{\mathrm{fl}}+2)\exp(-\varepsilon v)
               \le\frac1{16}(1-L^{-\varepsilon/2}).
 \end{gathered}
\end{equation}
In addition
$0<\varepsilon\le1/200$,
$c_{\mathrm{fl}}\ge1$,
$\bar e_0^\varepsilon\le1/4$,
$c_{\mathrm{fl}}L^{-3/8+6/200}\le1/16$, and
$L^{-5/8}\le1/4$.
The construction below applies whenever
$0<e\sqrt h\le\bar e_1$.

The physical top side and reference side are
\begin{equation}\label{eq:rgfl-reference-sides}
 R_{\mathrm{top}}=hMr_K,\qquad
 R^*=R_{\mathrm{top}}/L=hL^{T^*(e)},\qquad
 T^*(e)=\log_L(Mr_K)-1.
\end{equation}
They are independent of $N$, since
$v_K=-\log(e\sqrt h)$.
The admissible physical sides for the flow are
\[
 \mathcal T(e)=\{hL^T:T\in\N\cup\{0\},\ hL^T\le R_{\mathrm{top}}\}.
\]
For $R\in\mathcal T(e)$ put
$\ell_j(R)=R/a_j$,
$T_j(R)=\eta_j\Z^3/\ell_j(R)\Z^3$ and
$\mathfrak n_j=\log_L(\ell_j/(Mr_j))$.
The allowed unit sides are
$\mathrm{Cl}=\{\ell:\ell/(LM)\text{ is an odd integer at least }13\}$.
Their occurrence at every required scale is part
of the following arithmetic argument.

\begin{proposition}\label{prop:rgfl-scales}
For each $R\in\mathcal T(e)$ there is a finite integer
$N_{\min}(R,e)\ge0$ such that, when $N\ge N_{\min}(R,e)$,
the first level with $\mathfrak n_j=0$ exists and
satisfies $1\le k_*(R)\le K$.
Its physical quantities $a_*,e_*,\mu_*,\ell_*,r_*$
are independent of $N$ in this range.
For $j\le k_*(R)$,
\begin{equation}\label{eq:rgfl-scale-ranges}
 \begin{gathered}
 \ell_j\in\mathrm{Cl},\qquad
 \ell_j\ge Mr_j,\qquad r_j\ge r_*>31L+16,\\
 v_j\ge v_*\ge v_K,\qquad
 0<e_j\le e_*\le e_K\le\bar e_1,\qquad
 0<\mu_j\le\mu_*\le\bar m h\le1 .
 \end{gathered}
\end{equation}
For $j<k_*(R)$ one has $\ell_j\ge LMr_j$,
$\ell_{j+1}=\ell_j/L\in\mathrm{Cl}$,
$e_{j+1}=\sqrt L e_j\le\bar e_0$ and
$0<\mu_j=\mu_{j+1}/L\le1/L$.
At the stop,
\begin{equation}\label{eq:rgfl-terminal-scales}
 \ell_*=Mr_*,\qquad
 Mv_*^r\le\ell_*<LMv_*^r,\qquad
 \mu_*^{-1}<L\ell_*,\qquad e_*^2\ell_*=e^2R.
\end{equation}
If $R\le R'$ in $\mathcal T(e)$, then
$N_{\min}(R',e)\le N_{\min}(R,e)$ and
$k_*(R)\le k_*(R')$ for a common cutoff above both stated thresholds.
\end{proposition}

\begin{proof}
Use the depth $s=K-j$. Its quantities are
\[
 \widehat v_s=v_K+\tfrac{s}{2}\log L,\qquad
 \widehat r_s=\min\{L^i:i\ge0,\ L^i\ge\widehat v_s^r\},\qquad
 \widehat n_s(R)=T+s-\log_L(M\widehat r_s).
\]
They do not depend on $N$. At depth $s=K-j$ they
equal $v_j,r_j,\mathfrak n_j$, while
$a_j=hL^{-s}$ and $\ell_j=L^{T+s}$.
Since $\widehat v_s\ge v_K$, the scalar condition gives
\[
 \widehat v_{s+1}^r
   =\widehat v_s^r\left(1+\frac{\log L}{2\widehat v_s}\right)^r
                    \le L\widehat v_s^r\le L\widehat r_s.
\]
Thus
$\widehat r_{s+1}\in\{\widehat r_s,L\widehat r_s\}$
and $\widehat n_{s+1}-\widehat n_s\in\{0,1\}$.
Also $\widehat n_0\le0$ because $R\le R_{\mathrm{top}}$.
Minimality of the rounded power gives
$\widehat r_s<L\widehat v_s^r$ and hence
\[
 \widehat n_s>T+s-1-\log_LM-r\log_L(v_K+\tfrac{s}{2}\log L)
                                                    \longrightarrow\infty .
\]
The integer sequence therefore has a largest zero.
Denote its index by $s_*(R,e)$.
It is nonpositive up to that index and at least $1$
afterwards.

The condition $\ell_0\ge LMr_0$ is
$\widehat n_K\ge1$, equivalently $K>s_*$.
Its exact threshold is
\begin{equation}\label{eq:rgfl-mesh-threshold}
 N_{\min}(R,e)=\max\{0,s_*(R,e)+1-\log_Lh\},
             \qquad k_*(R)=K-s_*(R,e).
\end{equation}
This gives $1\le k_*\le K$ and
$a_*=hL^{-s_*}$, proving cutoff independence of
the terminal quantities.
Increasing $T$ increases every $\widehat n_s$
by the same integer; thus it decreases $s_*$,
decreases $N_{\min}$ and increases $k_*$.

For $j\le k_*$, the depth satisfies $K-j\ge s_*$.
Monotonicity gives the asserted ranges of
$\mathfrak n_j,r_j,v_j,e_j$, and the running mass
increases with $j$.
Moreover
$\bar m h=L^{\lfloor\log_L(1/\bar m)\rfloor-\log_L(1/\bar m)}
\in(1/L,1]$.
Since $\mu_*\ell_*=\bar mR\ge\bar m h>1/L$,
we obtain $\mu_*^{-1}<L\ell_*$.
The rounded-power inequality at $v_*$ gives
the two bounds for $\ell_*$ in
\eqref{eq:rgfl-terminal-scales}; its last identity
uses $e_*^2=e^2a_*$.

Finally $r_*\ge v_*^r>31L+16>L$ implies
$r_*\ge L^2$, as it is a power of $L$.
At every required level,
$\ell_j/(LM)=r_jL^{\mathfrak n_j-1}$
is a positive power of the odd integer $L\ge13$.
This proves $\ell_j\in\mathrm{Cl}$.
Before the stop $\mathfrak n_j\ge1$, giving the
stronger side bound. The mass and charge bounds
for the step follow at $j+1$.
\end{proof}

On $R^*$ define the reference mesh threshold
\begin{equation}\label{eq:rgfl-reference-threshold}
                    N_{\mathrm{ref}}(e)=N_{\min}(R^*(e),e).
\end{equation}
The stop on this torus has depth $s_*\in\{1,2\}$.
Indeed, for $s=1,2,3$ the stronger inequality in
\eqref{eq:rgfl-scale-smallness} gives
$\widehat r_s\in\{r_K,Lr_K\}$.
On $R^*$,
$\widehat n_s=s-1-\log_L(\widehat r_s/r_K)$.
Its values at $s=0,1,2,3$ are respectively
$-1$, an element of $\{-1,0\}$, an element of
$\{0,1\}$, and an element of $\{1,2\}$.
The zero-step-or-one-step rule forces the last zero
to be $1$ or $2$. On $R_{\mathrm{top}}$ the sequence
is larger by one, so its last zero is $0$ or $1$.
Consequently
\begin{equation}\label{eq:rgfl-reference-stop}
 \begin{gathered}
 k_*(R^*)\in\{K-1,K-2\},\qquad
 a_*\in\{h/L,h/L^2\},\qquad r_*\in\{r_K,Lr_K\},\\
 k_{\mathrm{top}}\in\{K,K-1\},\qquad
 L^{-3}<\mu_*\le L^{-1}\quad\text{on }R^* .
 \end{gathered}
\end{equation}
All these physical terminal values are independent
of $N\ge N_{\mathrm{ref}}(e)$.

The scale proposition verifies every side, mass
and charge hypothesis of the step theorem
on $R_{\mathrm{top}}$ and on the smaller tori it covers with odd degree.
At each level use the top system of
\cref{subsec:rgcov-top-system}; on a smaller
torus use its projection from
\cref{prop:rgcov-projected-system}.
These are the actual mean-field, covariance,
local resolvent and photon families constructed
earlier, with the same common estimates.
The ratio $R_{\mathrm{top}}/R$ is a power of the
odd integer $L$, and $k_*(R)\le k_{\mathrm{top}}$.
Thus the required top families exist whenever
a smaller-volume step is used.
Furthermore
$\ell_{j+1}(R)\ge Mr_{j+1}>(31L+16)M$;
the next-side hypothesis of the cover comparison
holds at every such step.

